\section{Introduction}

The classical global-regularity problem for the three-dimensional incompressible Navier--Stokes equations sits inside a long PDE line running from Leray and Hopf weak solutions through the monographs of Ladyzhenskaya and Temam, with partial regularity and continuation criteria due to Caffarelli--Kohn--Nirenberg, Prodi, Serrin, Beale--Kato--Majda, Escauriaza--Seregin--Sver\'ak, and Koch--Tataru isolating adjacent regularity mechanisms \cite{Leray1934,Hopf1951,Ladyzhenskaya1969,Temam2001,CaffarelliKohnNirenberg1982,Prodi1959,Serrin1962,BealeKatoMajda1984,EscauriazaSereginSverak2003,KochTataru2001}.

We consider the classical three-dimensional incompressible Navier-Stokes system
\begin{equation}
\partial_t u + (u\cdot\nabla)u + \nabla p = \nu\Delta u,\qquad \nabla\cdot u = 0.
\end{equation}
The domain is \([0,T)\times\mathbb{R}^3\), and the initial data are smooth, compactly supported, and divergence-free. The theorem proved here stays on that classical Euclidean surface throughout. In particular, results obtained only for regularized equations are treated as diagnostic context unless they are recovered on the classical system without theorem-critical modification.

The paper is organized around three bridge propositions. The first is a scale barrier that prevents indefinite transfer into dangerous high frequencies. The second is a compactness upgrade strong enough to pass the nonlinear term on the classical equation. The third is a same-surface gradient bound that excludes derivative-scale breakdown. The main theorem is obtained by discharging those three bridges on one admissibility surface and then closing the resulting singularity route.

The organizational choice is therefore simple: state one full-claim theorem, prove each bridge on the classical equation itself, and keep title, abstract, theorem, proof, and conclusion aligned with that same claim level.
